\subsection{Duality in modules of finite length over a principal ideal domain}

In this section A denotes a principal ideal domain which is not a field (and hence has at least one irreducible element), and K the field of fractions of A. For every A-module M, put

\[
\mathrm{D} (\mathbf {M}) = \operatorname{Hom} _ {\mathbf {A}} (\mathbf {M}, \mathbf {K} / \mathbf {A});
\]

we know that D(M) is equipped with the structure of an A-module in a natural way, namely for every homomorphism \(u: M \to K/A\) and every \(\alpha \in A\) , the homomorphism \(\alpha u\) maps x to \(\alpha u(x) = u(\alpha x)\) . To every homomorphism \(f: M \to N\) of A-modules is associated with the homomorphism D(f): D(N) → D(M), where D(f)(v) = v ∘ f (II, p.~196). For \(x \in M\) and \(x' \in D(M)\) , we put \(\langle x, x' \rangle = x'(x) \in K/A\) ; then \((x, x') \mapsto \langle x, x' \rangle\) is an A-bilinear map from M × D(M) into K/A, called canonical.

If M and N are two A-modules, then to every A-bilinear map \(\varphi: M \times N \to K/A\) are associated the A-linear maps \(d_{\varphi}: N \to D(M)\) and \(s_{\varphi}: M \to D(N)\) , where \(d_{\varphi}(y)(x) = \varphi(x, y) = s_{\varphi}(x)(y)\) (II, p.~268, Cor. to Prop. 1). In particular the canonical A-bilinear map \(M \times D(M) \to K/A\) defines an A-linear map (also called canonical)

\[
c _ {\mathbf {M}}: \mathbf {M} \rightarrow \mathbf {D} (\mathbf {D} (\mathbf {M}))
\]

such that \(\langle x', c_{\mathbf{M}}(x) \rangle = \langle x, x' \rangle\) for \(x \in \mathbf{M}\) and \(x' \in \mathrm{D}(\mathbf{M})\) .

PROPOSITION 10. --- If M is an A-module of finite length, then D(M) is (in general non-naturally) isomorphic to M, and the canonical map \(c_{\mathbf{M}}: \mathbf{M} \to \mathbf{D}(\mathbf{D}(\mathbf{M}))\) is an isomorphism.

Using VII, p.~19, Th. 2 and II, p.~203, Cor. 1, we reduce to the case where M is cyclic. Thus we may suppose that M = A/tA with \(t \neq 0\) . Note that any homomorphism \(u: A/tA \to K/A\) is completely determined by the image \(\xi \in K/A\) under u of the class \(\varepsilon\) of 1 mod tA, and this element must satisfy the relation \(t\xi = 0\) ; conversely, for any \(\xi \in K/A\) such that \(t\xi = 0\) , there exists a homomorphism \(u: A/tA \to K/A\) such that \(u(\varepsilon) = \xi\) . It follows that D(M) is isomorphic to \(t^{-1}A/A\) , and since the homothety by t is a bijection on K, we also have D(M) isomorphic to A/tA, which proves the first assertion. This proves that M and D(D(M)) are isomorphic, so have the same length; on the other hand \(c_{M}\) is injective, for if \(y \in A\) is such that the relation \(tz \in A\) (for \(z \in K\) ) implies \(yz \in A\) , then taking \(z = t^{-1}\) we have \(y \in tA\) . It follows that the image \(c_{M}(M)\) is necessarily equal to D(D(M)).

COROLLARY. --- Let M, N be two A-modules of finite length, and let \(\varphi\) be an A-bilinear map from \(M \times N\) into K/A, such that: 1) the relation \(\varphi(x, y) = 0\) for all \(y \in N\) implies x = 0; and 2) the relation \(\varphi(x, y) = 0\) for all \(x \in M\) implies y = 0. Then the A-linear maps \(s_{\varphi}: M \to D(N)\) and \(d_{\varphi}: N \to D(M)\) associated to \(\varphi\) are isomorphisms.

Indeed the hypotheses on \(\varphi\) mean that \(s_{\varphi}\) and \(d_{\varphi}\) are injective and since \(\text{long}(D(N)) = \text{long}(N)\) and \(\text{long}(D(M)) = \text{long}(M)\) by Prop. 10, this implies that \(\text{long}(M) = \text{long}(N)\) , and consequently \(s_{\varphi}\) and \(d_{\varphi}\) are bijective.

PROPOSITION 11. --- If \(M' \xrightarrow{u} M \xrightarrow{v} M''\) is an exact sequence of A-modules of finite length, then the sequence \(D(M'') \xrightarrow{D(v)} D(M) \xrightarrow{D(u)} D(M')\) is exact \(^{1}\) .

Let us show first of all that given an exact sequence

\[
0 \rightarrow \mathbf {M} ^ {\prime} \rightarrow \mathbf {M} \rightarrow \mathbf {M} ^ {\prime \prime} \rightarrow 0\tag{1}
\]

the corresponding sequence

\[
0 \rightarrow D (M ^ {\prime \prime}) \rightarrow D (M) \rightarrow D (M ^ {\prime}) \rightarrow 0
\]

is exact; indeed we know that the sequence

\[
0 \rightarrow \mathrm{D} (\mathbf {M} ^ {\prime \prime}) \rightarrow \mathrm{D} (\mathbf {M}) \rightarrow \mathrm{D} (\mathbf {M} ^ {\prime})
\]

is exact (II, p.~227, Th. 1); on the other hand, it follows from (1) that

\[
\operatorname{long} (\mathbf {M}) = \operatorname{long} \left(\mathbf {M} ^ {\prime}\right) + \operatorname{long} \left(\mathbf {M} ^ {\prime \prime}\right)
\]

(II, p.~212, Prop. 16); by Prop. 10 we thus have

\[
\operatorname{long} \left(\mathrm{D} (\mathbf {M})\right) = \operatorname{long} \left(\mathrm{D} \left(\mathbf {M} ^ {\prime}\right)\right) + \operatorname{long} \left(\mathrm{D} \left(\mathbf {M} ^ {\prime \prime}\right)\right),
\]

in other words long(D(M')) = long(D(M)/D(M'')). Since D(M)/D(M'') is naturally identified with a submodule of D(M'), it must be equal to D(M'), which proves our assertion.

This implies immediately that if \(u: \mathbf{M}' \to \mathbf{M}\) is injective, then \(D(u): D(\mathbf{M}) \to D(\mathbf{M}')\) is surjective; the conclusion then follows from II, p.~199, remark 4.

For every A-module M, let \(\mathfrak{S}(M)\) denote the set of submodules of M. For every submodule N of M (resp. every submodule \(N'\) of D(M)), let \(N^0\) (resp. \(N'^0\) ) denote the submodule of D(M) (resp. M) consisting of those \(x' \in D(M)\) (resp. \(x \in M\) ) such that \(\langle y, x' \rangle = 0\) for all \(y \in N\) (resp. \(\langle x, y' \rangle = 0\) for all \(y' \in N'\) ).

PROPOSITION 12. --- Let M be an A-module of finite length. Then the map which sends every submodule N of M to \(N^{0}\) is a bijection from \(\mathfrak{S}(M)\) onto \(\mathfrak{S}(D(M))\) , and the inverse bijection sends every submodule \(N'\) of D(M) to the submodule \(N'^{0}\) of M; the module D(N) is naturally identified with D(M)/ \(N^{0}\) and D(M/N) with \(N^{0}\) . Moreover, we have

\[
\left(\mathrm{N} _ {1} + \mathrm{N} _ {2}\right) ^ {0} = \mathrm{N} _ {1} ^ {0} \cap \mathrm{N} _ {2} ^ {0}, \quad \left(\mathrm{N} _ {1} \cap \mathrm{N} _ {2}\right) ^ {0} = \mathrm{N} _ {1} ^ {0} + \mathrm{N} _ {2} ^ {0}\tag{2}
\]

for all submodules \(\mathbf{N}_1\) , \(\mathbf{N}_2\) of \(\mathbf{M}\) .

For each submodule N of M, there is an exact sequence

\[
0 \rightarrow \mathrm{N} \rightarrow \mathrm{M} \rightarrow \mathrm{M} / \mathrm{N} \rightarrow 0
\]

and hence (Prop. 11) an exact sequence

\[
0 \rightarrow D (M / N) \rightarrow D (M) \rightarrow D (N) \rightarrow 0
\]

and since the image of D(M/N) in D(M) is obviously \(N^{0}\) , we see (Prop. 10) that \(\text{long}(N^{0}) = \text{long}(M) - \text{long}(N)\) ; since M is identified with D(D(M)) by Prop. 10, we have similarly

\[
\operatorname{long} \left(\mathbf {N} ^ {0 0}\right) = \operatorname{long} (\mathbf {M}) - \operatorname{long} \left(\mathbf {N} ^ {0}\right) = \operatorname{long} (\mathbf {N});
\]

in addition it is clear that \(N \subset N^{00}\) , so \(N^{00} = N\) . Furthermore the first relation in (2) is obvious, and by applying it to the submodules \(N_{1}^{0}\) and \(N_{2}^{0}\) of D(M), we have \((\mathbf{N}_{1}^{0} + \mathbf{N}_{2}^{0})^{0} = \mathbf{N}_{1} \cap \mathbf{N}_{2}\) , so \(N_{1}^{0} + N_{2}^{0} = (\mathbf{N}_{1}^{0} + \mathbf{N}_{2}^{0})^{00} = (\mathbf{N}_{1} \cap \mathbf{N}_{2})^{0}\) . This completes the proof of the proposition.

Examples. --- 1) For A = Z, the Z-modules of finite length are precisely the finite abelian groups; then K = Q, so K/A = Q/Z. Then to define D(M), we sometimes take, instead of Q/Z, a Z-module isomorphic to it, such as (V, p.~79, Prop. 2) the group R of roots of unity (under multiplication) in an algebraically closed field of characteristic 0; we then put D(M) = Hom \(_{Z}\) (M, R). We leave the reader to rewrite the preceding results for this special case in the corresponding notation.

\begin{enumerate}
\def\labelenumi{\arabic{enumi})}
\setcounter{enumi}{1}
\tightlist
\item
  Let \(a\) be a nonzero element of A. The map \(x \mapsto x/a\) from A into K induces an isomorphism on quotient modules from \(\mathbf{A}/(a)\) onto the submodule \((\mathbf{K}/\mathbf{A})(a)\) of K/A consisting of elements annihilated by \(a\) . If M is an A-module annihilated by \(a\) , or equivalently an \(\mathbf{A}/(a)\) -module, then the A-module D(M) is identified with \(\operatorname{Hom}_{\mathbf{A}/(a)}(\mathbf{M}, \mathbf{A}/(a))\) . We leave the reader to rewrite the preceding results for this special case in the corresponding notation (cf.~V, p.~86).
\end{enumerate}

